\chapter{MIDI input}\label{ch:midi}

\RITMO{} plays and records notes from a MIDI keyboard. The input device is chosen in \menu{Tools $\rightarrow$ Options} (the MIDI
group); the button \menu{Toggle MIDI on/off} of the toolbar switches the input on and off while a device is set. The device is a port
of the system: on Linux the ports of the ALSA sequencer (a USB keyboard or a virtual port), on Windows WinMM, on macOS CoreMIDI.

\section{Options}

\begin{center}
\begin{tabular}{@{}lp{10.2cm}@{}}
\toprule
Option & Meaning \\ \midrule
MIDI IN device & the input device, or \emph{none} \\
Touch response & the key velocity becomes the volume of the note: \texttt{Atari Volume Offset + velocity / 8}, limited to 1--15. Without it the
  current volume (as for the typed notes) is used \\
Atari volume offset & added to the volume of the velocity, from 0 to 15 \\
Record note off & a released key (a note off, or a note on with velocity 0) deletes the note it played and writes volume 0. Without it releases are
  ignored \\
\bottomrule
\end{tabular}
\end{center}

\section{The MIDI channels}

\begin{center}
\begin{tabular}{@{}lp{10.5cm}@{}}
\toprule
Channel & Function \\ \midrule
1 & Records notes at the cursor, like the note keys: the MIDI note 36 is the lowest note of the tracker (\texttt{C-1}), the note is written with the
  active instrument, the cursor moves down by the note spacing, and the note is played. While the song plays with follow mode, a note in the
  first half of a line is written on the next line (quantized). Outside the tracks area, in a jam mode or with \key{SHIFT} or \key{CONTROL} held down, the note is only
  played. Recording needs the window to have the focus \\
2--9 & Live play on the channels \texttt{L1}--\texttt{L4} and \texttt{R1}--\texttt{R4} (6--9 repeat \texttt{L1}--\texttt{L4} on a mono song): never
  recorded, and independent of the focus. The velocity is the volume (velocity / 8); a program change sets the instrument of the channel;
  the controller 123 (all notes off) stops its note and the controller 121 (reset all controllers) restarts the player routine \\
11--15 & a program change chooses the active instrument \\
10 and 16 & an experimental assignment of controllers, kept from the original program, with the \texttt{POKEY explorer} (chapter~\ref{ch:explorer}) on
  channel 16 \\
\bottomrule
\end{tabular}
\end{center}

A system reset message (\texttt{FF}) restarts the player routine and forgets the notes held on the channels 2 to 16. The implementation chart of
the MIDI messages is in appendix~\ref{app:midi}.

\section{Scripts}

A script can send MIDI messages to the program with the \texttt{midi} command, which is how the recording can be tested without a keyboard
(chapter~\ref{ch:scripting}).
